Given a training function $f$, \replay{} enables us to compute metagradients
$\nabla f(\mathbf{z})$ for any setup $\mathbf{z}$. Can we immediately
use these metagradients to
optimize model training setups?
The answer is (generally) no:
we find that applying \replay{} to a function $f$ representing a standard model
training and evaluation routine yields metagradients that are often
$\pm\infty$-valued and generally unhelpful for optimization.
Indeed, previous work has observed similar issues optimizing
over even (very) small-scale training
\citep{bengio1994learning,pearlmutter1996investigation,maclaurin2015gradient}.

In this section, we show that an underlying source of the issue is the
{landscape} of the metaparameter optimization problem.
We then present a framework for modifying standard learning
algorithms
to admit useful metagradients, i.e., to be {\em metasmooth}.
To use a familiar analogy:
just as residual connections and improved initialization schemes can improve
optimization in standard deep learning algorithms, our framework introduces an
analogous set of modifications to enable optimization with metagradients.

\subsection{The metaparameter optimization landscape}
\label{sec:smoothness}
We first review the notion of smoothness from optimization theory,
and then adapt it to the setting of metagradients.
The resulting {\em metasmoothness} metric allows us to quantify
(and later, improve) the amenability
of the metaparameter optimization problem to gradient-based methods.

\paragraph{Smoothness.}
In optimization theory,
the basic property of a function that controls how effectively it can
be optimized
with first-order methods is {\em smoothness}.
Specifically, a function $f(\mathbf{z})$ is $\beta$-smooth at a point
$\mathbf{z}$
if its gradient $\nabla f$ satisfies the property that
\begin{equation}
  \label{eq:smoothness}
  \|\nabla f(\mathbf{z}) - \nabla f(\mathbf{z}')\| \leq \beta \cdot
  \|\mathbf{z} - \mathbf{z}'\| \qquad \text{ for all $\mathbf{z}'$},
\end{equation}
or in other words, if its gradient does not change too quickly around
$\mathbf{z}$.
To motivate this definition:
if a function $f$ is
$\beta$-smooth at $\mathbf{z}$, then a step of gradient descent with step size
$\nicefrac{1}{\beta}$ will successfully decrease the value of the function:
\begin{align*}
  f\left(\mathbf{z} - \frac{1}{\beta} \nabla f(\mathbf{z})\right)
  \leq f(\mathbf{z}) - \frac{1}{2\beta}\|\nabla f(\mathbf{z})\|^2.
\end{align*}
This guarantee makes $\beta$-smoothness
a good measure of gradient utility.

\paragraph{Metasmoothness.}
There are two main challenges in adapting the
smoothness property to the
metagradient setting.
First, evaluating \eqref{eq:smoothness} requires a search over all
possible $\mathbf{z}'$,
which is infeasible.
Second, even if we could exactly evaluate the left-hand side of
\eqref{eq:smoothness},
it would be difficult to disentangle non-smoothness of the training function $f$
from potential error in metagradient computation (e.g., a numerically unstable
operation in \replay{}).

To sidestep these issues, we propose a metric called {\em metasmoothness},
given in Definition \ref{def:meta_smoothness}.
Metasmoothness is cheap to compute---requiring only three evaluations of the
training function---and does not rely on metagradient computation.
For the remainder of this section, we fix a small constant $h > 0$,
and define the corresponding
finite-differences estimator of the directional derivative $\Delta_f$ as
\begin{align*}
  \Delta_f(\mathbf{z}; \mathbf{v}) := \frac{f(\mathbf{z} +
  h\mathbf{v}) - f(\mathbf{z})}{h}.
\end{align*}

\begin{definition}[Metasmoothness of $f$ at $\mathbf{z}$ towards $\mathbf{v}$]
  \label{def:meta_smoothness}
  Consider a training function $f$ mapping metaparameters $\mathbf{z}
  \in \mathbb{R}^n$
  to model output $f(\mathbf{z}) \in \mathbb{R}$.
  Given a metaparameter $\mathbf{z}$ and
  a vector $\mathbf{v} \in \mathbb{R}^n$,
  the metasmoothness of $f$ at $\mathbf{z}$ towards $\mathbf{v}$ is given by
  \begin{align}
    \label{eq:meta_smoothness_approx}
    S_{h, \mathbf{v}}(f; \mathbf{z})
    &\coloneqq
    \bigg|\frac{\Delta_f(\mathbf{z} + h \mathbf{v}) -
    \Delta_f(\mathbf{z})}{h} \bigg|.
  \end{align}
\end{definition}
\noindent
Definition \ref{def:meta_smoothness} measures the rate of change of
the derivative of $f(\mathbf{z})$ in the direction of a given vector
$\mathbf{v}$,
and is therefore related to $\beta$-smoothness in that:
\begin{itemize}
  \item[(a)] If $f$ is $\beta$-smooth at $\mathbf{z}$, then $S_{h,
    \mathbf{v}}(f; \mathbf{z}) \leq \beta$ for any $(h, \mathbf{v})$
    (so Definition \ref{def:meta_smoothness} is {\em necessary} for smoothness).
  \item[(b)] If $\lim_{h \to 0} S_{h, \mathbf{v}}(f; \mathbf{z}) \leq
    \beta$ for all $\mathbf{z} \in \mathbb{R}^n$ and $\mathbf{v} \in
    \mathbb{S}^{n-1}$,
    then $f$ is $\beta$-smooth everywhere
    (so a global version of Definition \ref{def:meta_smoothness} is
    {\em sufficient} for smoothness).
\end{itemize}

\paragraph{Empirical metasmoothness.}
Definition \ref{def:meta_smoothness} lets us measure the
meta-smoothness of a
training function $f$ at a particular metaparameter $\mathbf{z}$
(towards a direction $\mathbf{v}$).
This definition, however, has two shortcomings.
First, recall that the training function $f$ is a composition of a learning
algorithm $\mathcal{A}$ and an output function $\phi$,
so the smoothness of $f$ depends on that of both $\mathcal{A}$
and $\phi$ (in particular,
  $\nicefrac{\partial f}{\partial \mathbf{z}} = \nicefrac{\partial
  \phi}{\partial
\mathcal{A}}\cdot \nicefrac{\partial \mathcal{A}}{\partial \mathbf{z}}$).
Since the output
function $\phi$ might be unknown ahead of time, we are most interested in
measuring the {\em overall} metasmoothness of a learning algorithm
$\mathcal{A}$.
Second, while the result of \eqref{eq:meta_smoothness_approx} does have a
concrete basis in optimization theory, it may not be easy to
interpret in practice
(e.g., what does $S = 200$ mean?).
We address both issues simultaneously by
(a) proposing an interpretable ``binarized''
version of Definition \ref{def:meta_smoothness}, and
(b) studying metasmoothness in the space of model parameters $\theta$,
instead of the output space.
\begin{definition}[Empirical metasmoothness of $\mathcal{A}$]
  \label{def:avg_meta_smoothness}
  Let $\mathcal{A}$ be a learning algorithm which maps metaparameters
  $\mathbf{z} \in \mathbb{R}^n$
  to model parameters $\theta \in \mathbb{R}^d$,
  let $\mathbf{z}$ be a metaparameter vector,
  and let $\mathbf{v}$ be a given direction.
  Let $\mathbf{d} \in \mathbb{R}^d$ be the
  per-coordinate variation in $\theta$, i.e.,
  \begin{align*}
    \mathbf{d} = |\mathcal{A}(\mathbf{z} + 2h\mathbf{v}) -
    \mathcal{A}(\mathbf{z})|  %
  \end{align*}
  The empirical $(h, \mathbf{v})$-metasmoothness of $\mathcal{A}$ at
  $\mathbf{z}$ is given by
  \begin{align}
    \label{eq:avg_meta_smoothness}
    \widehat{S}_{h, \mathbf{v}}(\mathcal{A}; \mathbf{z})
    & =
    \mathrm{sign}(\Delta_{\mathcal{A}}(\mathbf{z};\mathbf{v}))^\top\cdot
    \mathrm{diag}\left(
      \frac{\mathbf{d}}{\|\mathbf{d}\|_1}
    \right)\cdot \mathrm{sign}(\Delta_{\mathcal{A}}(\mathbf{z} +
    h\mathbf{v};\mathbf{v})),
  \end{align}
  weights each parameter by its range.
\end{definition}
Intuitively, \eqref{eq:avg_meta_smoothness} measures the
agreement in sign between the (finite-difference approximation
of the) metagradient in the direction of $\mathbf{v}$ at $\mathbf{z}$
and at $\mathbf{z} + h\mathbf{v}$,
averaged across parameter coordinates and weighted by
the variation in each coordinate.
Taking a weighted average of sign agreements ensures that
$\smash{\widehat{S} \in [-1, 1]}$ (making it easier to interpret than
Definition \ref{def:meta_smoothness}).
The $\smash{\mathrm{diag}(\nicefrac{\mathbf{d}}{\|\mathbf{d}\|_1})}$ term
weights each agreement proportionally to the scale of the
corresponding parameter change (downweighting, e.g., coordinates $i$ that
are essentially constant).
Finally, observe that Definition \ref{def:avg_meta_smoothness} is efficient
to compute in practice: it requires only three calls to the learning
algorithm $\smash{\mathcal{A}}$.


\begin{remark}[]
  Ideally, recalling the smoothness definition \eqref{eq:smoothness},
  we would evaluate metasmoothness in all possible directions
  $\mathbf{v}$ and all points $\mathbf{z}$. Empirically, we find in
  the sequel (Section \ref{sec:improving})
  that this single-direction approximation at a single point $\mathbf{z}$
  still yields a useful estimate of metasmoothness
  (e.g., one that correlates with metagradient utility).
\end{remark}

\begin{figure}
  \centering
  \begin{subfigure}[t]{0.61\textwidth}
    \captionsetup{labelformat=empty}%
    \pgfplotsset{colormap/Set2}
\pgfplotsset{grid style={dashed,gray}}

\begin{tikzpicture}
    \begin{axis}[
      xlabel={Accuracy (\%)},
      ylabel={Smoothness},
      legend style={at={(1.1,1)}, anchor=north west, inner sep={6pt}},
      legend cell align={left},
      legend columns={3},
      height={6.8cm},
      width={6cm},
      colormap name={Set2},
      grid={both}
    ]

    \pgfkeys{
        /markarray/0/.initial=*,
        /markarray/1/.initial=triangle*,
        /markarray/2/.initial=square*,
        /markarray/3/.initial=o,
        /markarray/4/.initial=square*,
        /markarray/5/.initial=square
    } 
    \addplot [
      scatter,
      only marks,
      forget plot,
      visualization depends on=\thisrow{bs} \as \bsize,
      visualization depends on=\thisrow{final_scale} \as \fscale,
      visualization depends on=\thisrow{batchnorm_before_act} \as \bnact,
      visualization depends on=\thisrow{width_multiplier} \as \width,
      scatter/@pre marker code/.code={
        \pgfmathparse{\width == 1.0 ? 0 : (\width == 1.5 ? 1 : 2)}
        \edef\currentmark{\pgfkeysvalueof{/markarray/\pgfmathresult}}
        \pgfmathsetmacro{\mycolor}{\fscale == 0.125 ? 0 : (\fscale == 0.5 ? 128 : 255)}
        \pgfmathsetmacro{\myopacity}{\bnact == 0 ? 0.2 : 1}
        \pgfplotsset{mark=\currentmark}
        \def\markopts{mark size={(\bsize == 250 ? 2 : (\bsize == 500 ? 3 : 4))},%
        color of colormap=\mycolor,
        draw=black,
        opacity=\myopacity}
        \expandafter\scope\expandafter[\markopts]
      },
      scatter/@post marker code/.code={
        \endscope
      }, 
    ]
    table[
      x expr={\thisrow{avg_acc}*100},
      y=smoothness,
      col sep=comma,
    ]{pgffigs/smoothness_vs_acc/smoothness_vs_acc.csv};
    
    \addlegendimage{only marks, mark=none, draw=none, mark options={fill = white, draw = white}}%
    \addlegendentry{\makebox[0pt][l]{\textbf{Final scale}}}
    \addlegendimage{empty legend}
    \addlegendentry{}
    \addlegendimage{empty legend}
    \addlegendentry{}
    \addlegendimage{only marks, mark=*, mark options={color of colormap=0}}%
    \addlegendentry{0.125}
    \addlegendimage{only marks, mark=*, mark options={color of colormap=128}}%
    \addlegendentry{0.5}
    \addlegendimage{only marks, mark=*, mark options={color of colormap=255}}%
    \addlegendentry{1.0}

    \addlegendimage{only marks, mark=none, draw=none, mark options={fill = white, draw = white}}%
    \addlegendentry{$\phantom{\bigg|}$\makebox[0pt][l]{\textbf{Batch size}}}
    \addlegendimage{empty legend}
    \addlegendentry{}
    \addlegendimage{empty legend}
    \addlegendentry{}
    \addlegendimage{only marks, mark=*, mark size=2, mark options={fill=black}}%
    \addlegendentry{250}
    \addlegendimage{only marks, mark=*, mark size=3, mark options={fill=black}}%
    \addlegendentry{500}
    \addlegendimage{only marks, mark=*, mark size=4, mark options={fill=black}}%
    \addlegendentry{1000}

    \addlegendimage{only marks, mark=none, draw=none, mark options={fill = white, draw = white}}%
    \addlegendentry{$\phantom{\bigg|}$\makebox[0pt][l]{\textbf{Width}}}
    \addlegendimage{empty legend}
    \addlegendentry{}
    \addlegendimage{empty legend}
    \addlegendentry{}
    \addlegendimage{only marks, mark=*, mark size=3, mark options={fill=black}}%
    \addlegendentry{1.0}
    \addlegendimage{only marks, mark=triangle*, mark size=3, mark options={fill=black}}%
    \addlegendentry{1.5}
    \addlegendimage{only marks, mark=square*, mark size=3, mark options={fill=black}}%
    \addlegendentry{2.0}

    \addlegendimage{only marks, mark=none, draw=none, mark options={fill = white, draw = white}}%
    \addlegendentry{$\phantom{\bigg|}$\makebox[0pt][l]{\textbf{BN Placement}}}
    \addlegendimage{empty legend}
    \addlegendentry{}
    \addlegendimage{empty legend}
    \addlegendentry{}

    \addlegendimage{only marks, mark=*, mark size=3, opacity=1, mark options={fill=black}}%
    \addlegendentry{\makebox[0pt][l]{Before activation}}
    \addlegendimage{empty legend}
    \addlegendentry{}
    \addlegendimage{empty legend}
    \addlegendentry{}

    \addlegendimage{only marks, mark=*, mark size=3, opacity=0.2, mark options={fill=black}}%
    \addlegendentry{\makebox[0pt][l]{After activation}}
    \addlegendimage{empty legend}
    \addlegendentry{}
    \addlegendimage{empty legend}
    \addlegendentry{}

    \end{axis}
\end{tikzpicture}

    \caption{(a)\quad\quad\quad\quad\quad\quad\quad\quad\quad\quad\quad\quad\quad}
  \end{subfigure}
  \begin{subfigure}[t]{0.37\textwidth}
    \captionsetup{labelformat=empty}%
    \raisebox{0.0135\height}{\pgfplotsset{colormap/Set2}
\pgfplotsset{scaled y ticks=false}
\pgfplotsset{grid style={dashed,gray}}

\begin{tikzpicture}
    \begin{axis}[
      xlabel={Smoothness},
	  ylabel={$\Delta$ Loss (Optimization Objective)},
      height={6.8cm},
      width={6cm},
	  ytick pos=right,
	  y tick label style={
        /pgf/number format/fixed,     %
        /pgf/number format/precision=2,  
        /pgf/number format/zerofill   %
      },
      scaled y ticks=false,
      colormap name={Set2},
      grid={both}
    ]

    \pgfkeys{
        /markarray/0/.initial=*,
        /markarray/1/.initial=triangle*,
        /markarray/2/.initial=square*,
        /markarray/3/.initial=o,
        /markarray/4/.initial=square*,
        /markarray/5/.initial=square
    } 
    \addplot [
      scatter,
      only marks,
      forget plot,
      visualization depends on=\thisrow{bs} \as \bsize,
      visualization depends on=\thisrow{final_scale} \as \fscale,
      visualization depends on=\thisrow{batchnorm_before_act} \as \bnact,
      visualization depends on=\thisrow{width_multiplier} \as \width,
      scatter/@pre marker code/.code={
        \pgfmathparse{\width == 1.0 ? 0 : (\width == 1.5 ? 1 : 2)}
        \edef\currentmark{\pgfkeysvalueof{/markarray/\pgfmathresult}}
        \pgfmathsetmacro{\mycolor}{\fscale == 0.125 ? 0 : (\fscale == 0.5 ? 128 : 255)}
        \pgfmathsetmacro{\myopacity}{\bnact == 0 ? 0.2 : 1}
        \pgfplotsset{mark=\currentmark}
        \def\markopts{mark size={(\bsize == 250 ? 2 : (\bsize == 500 ? 3 : 4))},%
        color of colormap=\mycolor,
        draw=black,
        opacity=\myopacity}
        \expandafter\scope\expandafter[\markopts]
      },
      scatter/@post marker code/.code={
        \endscope
      }, 
    ]
    table[
      x={smoothness},
      y={val_acc_x},
      col sep=comma,
    ]{pgffigs/smoothness_vs_opt/smoothness_vs_opt.csv};
    \end{axis}
\end{tikzpicture}
}
    \caption{(b)\quad\quad\quad\quad}
  \end{subfigure}
  \caption{
    \textit{(a)} For a variety of training configurations of a ResNet-9
    model, we plot metasmoothness (Def. \ref{def:avg_meta_smoothness})
    against test accuracy. Strategies such as increasing width, placing
    batch normalization before activations, and scaling down network
    outputs consistently improve  metasmoothness, at a minor cost to
    accuracy.
    \textit{(b)} Smoother training configurations can be optimized via
    metagradients more effectively. Here, as in Section
    \ref{sec:poisoning}, we use metagradients to gradient ascend on
    validation loss.
  }
  \label{fig:smoothness_vs_accuracy}
\end{figure}


\subsection{Estimating and improving metasmoothness}
\label{sec:improving}
Having established a method for quantifying metasmoothness, we turn to the
practical question: how can we design learning algorithms that are amenable to
metagradient optimization? To answer this question,
we introduce a
straightforward
framework: given a learning algorithm,
explore a fixed menu of possible modifications to the
training setup, and choose the combination that maximizes empirical
metasmoothness.
In practice, we find that this framework allows us to slightly modify learning
algorithms in a way that makes them amenable to first-order methods.

As a case study, we
study the task of training ResNet-9 on the CIFAR-10 dataset
\citep{krizhevsky2009learning}.
We let the metaparameters $\mathbf{z}$ be
a perturbation to the pixels of 1000 random training images
(so $\smash{\mathbf{z} \in \mathbb{R}^{1000 \times 32 \times 32 \times 3}}$).
We estimate
the empirical metasmoothness of different learning algorithms $\mathcal{A}$
at $\mathbf{z} = \bm{0}$ using Definition \ref{def:avg_meta_smoothness}.
Concretely, we proceed as follows for each learning algorithm $\mathcal{A}$:
\begin{enumerate}
    \item Let $\mathbf{z}_0 = \bm{0}$ be the metaparameter corresponding to
    the original dataset.
    \item Sample a random perturbation vector $\mathbf{v} \sim \mathcal{N}(0, 1)$.
    \item Compute the empirical metasmoothness \eqref{eq:avg_meta_smoothness}, i.e.,
    \begin{enumerate}
        \item Let $\theta_0 := \mathcal{A}(\mathbf{z}_0)$,
        $\theta_h := \mathcal{A}(\mathbf{z}_0 + h \cdot \mathbf{v})$, and
        $\theta_{2h} := \mathcal{A}(\mathbf{z}_0 + 2h \cdot \mathbf{v})$ be the model parameters
        that result from training with training dataset perturbations
        $\mathbf{z}_0$, $\mathbf{z}_0 + h \mathbf{v}$,
        and $\mathbf{z}_0 + 2h \mathbf{v}$, respectively.
        \item Compute the approximate derivatives
        \[
            \Delta_{\mathcal{A}}(\mathbf{z}_0; \mathbf{v}) = \left(\theta_h - \theta_0\right)/h, \quad
            \Delta_{\mathcal{A}}(\mathbf{z}_0 + h \mathbf{v}; \mathbf{v}) = \left(\theta_{2h} - \theta_h\right)/h.
        \]
        \item Compute the weighting vector $\mathbf{d}  = |\theta_{2h} - \theta_0|$, and compute the average metasmoothness \eqref{eq:avg_meta_smoothness}, i.e.,
        \[
            \widehat{S}_{h, \mathbf{v}}(\mathcal{A}; z_0) =
            \textrm{sign}(\Delta_{\mathcal{A}}(\mathbf{z}_0 + h \mathbf{v};\mathbf{v}))^\top
            \cdot \textrm{diag}\left(\frac{\mathbf{d}}{\|\mathbf{d}\|_1}\right)
            \cdot \textrm{sign}(\Delta_{\mathcal{A}}(\mathbf{z}_0;\mathbf{v})).
        \]
    \end{enumerate}
\end{enumerate}



\begin{figure}[h]
  \centering
  \newcommand{\plotnumber}{34}
  \newcommand{\zmin}{0}
  \newcommand{\zmax}{12.5}
  \newcommand{\zdist}{3}
  \begin{tikzpicture}
    \centering
    \pgfplotsset{
        view={45}{30},
        mesh/ordering=y varies,
        table/col sep=comma,
        table/header=true,
    }

    \pgfplotsset{
        colormap={crest}{
            rgb=(0.647,0.803,0.566) %
            rgb=(0.500,0.737,0.568) %
            rgb=(0.357,0.651,0.565) %
            rgb=(0.247,0.586,0.562) %
            rgb=(0.154,0.471,0.548) %
            rgb=(0.113,0.355,0.518) %
            rgb=(0.132,0.289,0.486) %
            rgb=(0.173,0.190,0.445) %
        }
    }

    \pgfplotstableread[col sep=comma,header=true]{pgffigs/smoothness_3d/data/nonsmooth_3dplot_\plotnumber.csv}\nonsmoothdata
    \pgfplotstableread[col sep=comma,header=true]{pgffigs/smoothness_3d/data/smooth_3dplot_\plotnumber.csv}\smoothdata

    \begin{groupplot}[
        group style={
            group size=2 by 1,
            horizontal sep=0.1\textwidth,
            vertical sep=0.1\textwidth,
        },
        width=0.45\textwidth,
        height=0.35\textwidth,
        colormap name=viridis,
        xlabel=$h_1$,
        ylabel=$h_2$,
        zlabel=$f(h_1\mathbf{v}_1 + h_2\mathbf{v}_2)$,
        xlabel shift=-10pt,
        ylabel shift=-10pt,
        zlabel shift=-5pt,
        title style={yshift=-5pt},
        scaled ticks=true,
        xmin=0.001,
        ymin=0.001,
        zmin=\zmin,
        zmax=\zmax,
        grid=both,
        grid style={line width=0.5pt, draw=gray!30, dashed},
        xtick distance=0.25,
        ytick distance=0.25,
        ztick distance=\zdist,
        extra x ticks={0},
        extra y ticks={0},
        extra z ticks={0},
        extra tick style={grid=none, tick style={draw=none}, major tick length=0pt},
        extra x tick labels={},
        extra y tick labels={},
        extra z tick labels={}
    ]
        \nextgroupplot[title={Non-smooth (Example \#\plotnumber)}]
            \addplot3[
                surf,
                shader=interp,
                mesh/rows=21,
                mesh/cols=21
            ] table[x=X, y=Y, z=Z] {\nonsmoothdata};
            
        \nextgroupplot[title={Smooth (Example \#\plotnumber)}]
            \addplot3[
                surf,
                shader=interp,
                mesh/rows=21,
                mesh/cols=21
            ] table[x=X, y=Y, z=Z] {\smoothdata};
    \end{groupplot}
\end{tikzpicture}


  \vspace{.25cm} %

  \renewcommand{\plotnumber}{811}
  \renewcommand{\zmin}{-1}
  \renewcommand{\zmax}{3}
  \renewcommand{\zdist}{1}
  \begin{tikzpicture}
    \centering
    \pgfplotsset{
        view={45}{30},
        mesh/ordering=y varies,
        table/col sep=comma,
        table/header=true,
    }

    \pgfplotsset{
        colormap={crest}{
            rgb=(0.647,0.803,0.566) %
            rgb=(0.500,0.737,0.568) %
            rgb=(0.357,0.651,0.565) %
            rgb=(0.247,0.586,0.562) %
            rgb=(0.154,0.471,0.548) %
            rgb=(0.113,0.355,0.518) %
            rgb=(0.132,0.289,0.486) %
            rgb=(0.173,0.190,0.445) %
        }
    }

    \pgfplotstableread[col sep=comma,header=true]{pgffigs/smoothness_3d/data/nonsmooth_3dplot_\plotnumber.csv}\nonsmoothdata
    \pgfplotstableread[col sep=comma,header=true]{pgffigs/smoothness_3d/data/smooth_3dplot_\plotnumber.csv}\smoothdata

    \begin{groupplot}[
        group style={
            group size=2 by 1,
            horizontal sep=0.1\textwidth,
            vertical sep=0.1\textwidth,
        },
        width=0.45\textwidth,
        height=0.35\textwidth,
        colormap name=viridis,
        xlabel=$h_1$,
        ylabel=$h_2$,
        zlabel=$f(h_1\mathbf{v}_1 + h_2\mathbf{v}_2)$,
        xlabel shift=-10pt,
        ylabel shift=-10pt,
        zlabel shift=-5pt,
        title style={yshift=-5pt},
        scaled ticks=true,
        xmin=0.001,
        ymin=0.001,
        zmin=\zmin,
        zmax=\zmax,
        grid=both,
        grid style={line width=0.5pt, draw=gray!30, dashed},
        xtick distance=0.25,
        ytick distance=0.25,
        ztick distance=\zdist,
        extra x ticks={0},
        extra y ticks={0},
        extra z ticks={0},
        extra tick style={grid=none, tick style={draw=none}, major tick length=0pt},
        extra x tick labels={},
        extra y tick labels={},
        extra z tick labels={}
    ]
        \nextgroupplot[title={Non-smooth (Example \#\plotnumber)}]
            \addplot3[
                surf,
                shader=interp,
                mesh/rows=21,
                mesh/cols=21
            ] table[x=X, y=Y, z=Z] {\nonsmoothdata};
            
        \nextgroupplot[title={Smooth (Example \#\plotnumber)}]
            \addplot3[
                surf,
                shader=interp,
                mesh/rows=21,
                mesh/cols=21
            ] table[x=X, y=Y, z=Z] {\smoothdata};
    \end{groupplot}
\end{tikzpicture}

  \caption{
    The effect of metasmoothness on the optimization landscape.
    Each plot above visualizes the loss landscape of a (deterministic)
    learning algorithm $\mathcal{A}$, with the $x$- and $y$-axes representing
    additive perturbations to 1000 examples in the training set
    and the $z$-axis representing the resulting model's loss
    on the test example given in the title.
    In each row, the left plot is a non-smooth algorithm,
    and the right plot is a smooth algorithm
    (as per Definition \ref{def:avg_meta_smoothness})
    evaluated on the same  example.
    Overall, empirical metasmoothness seems to strongly correlate with
    qualitative landscape smoothness. See Figure \ref{fig:more_landscapes}
    for more examples.
  }
  \label{fig:smoothness_3d}
\end{figure}

\paragraph{Metasmooth learning algorithms.}
We apply the procedure above to estimate the metasmoothness of learning
algorithms induced by different design choices
(batch size, network width, BatchNorm placement, gradient scaling),
and report the results in Figure \ref{fig:smoothness_vs_accuracy}
(left). On one hand,
``standard'' learning algorithms (i.e., those designed without metasmoothness in
mind) are not metasmooth.
On the other hand, our investigation reveals central factors driving
metasmoothness.
In addition to ``standard''
hyperparameters such as batch size and network width playing a role, we find
that placing Batch Normalization layers {\em prior} to nonlinearities (instead
of after) and scaling the final layer output are both crucial to metasmoothness.
Note that the modifications we consider above are not exhaustive---see
Appendix \ref{app:poisoning} for the full training setup.

Finally, in Figure~\ref{fig:smoothness_3d}, we plot
the optimization landscape of both metasmooth (right) and
non-metasmooth (left) models.
We find that the landscapes of metasmooth models are much smoother
and---qualitatively---more straightforward to optimize.

\paragraph{Metasmoothness/performance tradeoffs?}
Figure \ref{fig:smoothness_vs_accuracy} (left) relates
metasmoothness to model accuracy for the considered learning algorithms.
While there is no clear trend, the top-performing
learning algorithms are not always metasmooth.
However, the trade-off is not too severe: the most metasmooth
algorithms still achieve near-optimal accuracy. Furthermore, it is possible
that with additional searching we could identify even more accurate metasmooth
models.
Taken together with our previous experiment, our results suggest that
jointly searching
over metasmoothness and model accuracy is a general recipe for designing
learning algorithms that are both performant and metasmooth.
Finally, as we discuss in Section \ref{sec:discussion},
a fruitful avenue for future work may be to
design metasmooth learning algorithms directly,
i.e., without relying on stability heuristics or grid search.


\paragraph{Does metasmoothness aid downstream optimization?}
Recall that our motivation for studying metasmoothness is to develop learning
algorithms that we can optimize the metaparameters of via metagradients (using
first-order methods). We started with the notion of $\beta$-smoothness from
optimization theory, and we adapted it to the setting of metagradients by making
a series of approximations and modifications. The final question we address is:
does our final notion of metasmoothness actually predict the utility of
metagradients for optimization? Figure~\ref{fig:smoothness_vs_accuracy}
(right) demonstrates
that metasmoothness strongly predicts our ability to optimize the
metaparameters of a
given learning algorithm. We use metagradients (computed by
\replay{}) to gradient
ascend on validation loss with respect to the metaparameters $\mathbf{z}$, and
measure the change in model loss.


